\section{Formal Properties}
\label{sec:properties}

This section delivers Contributions~2,~3, and~4. We prove a semantically
strong Parser Independence theorem, derive the layer dependency rule from
its contrapositive, formalize attacker-durability, characterize the distinct
durability profiles of \src{Q} and \src{C}, and give a schema-conscious
definition of hash-pivot.

\subsection{Parser Independence}

The theorem we want is not merely a type-checking observation. We need a
\emph{semantic} claim: that two distinct navigation paths into the
``same'' artifact yield the \emph{same} record under a fixed parser. We
make this precise with the following theorem, which has two halves -- a
syntactic well-typedness half, and a semantic equality half.

\begin{theorem}[Parser Independence]
\label{thm:parser-indep}
Let $p : \Bstar \to R$ be a deterministic, side-effect-free parser, and let
$S = (A, f_{\mathrm{nav}}, p)$ and $S' = (A', f'_{\mathrm{nav}}, p)$ be two
evidence sources with possibly distinct navigation primitives $\pi$ and
$\pi'$. Then:

\smallskip
\noindent\textbf{(I) Type composition.}
Both compositions $p \circ f_{\mathrm{nav}}: A \rightharpoonup R$ and $p
\circ f'_{\mathrm{nav}}: A' \rightharpoonup R$ are well-typed and yield
records in $R$.

\smallskip
\noindent\textbf{(II) Semantic invariance.}
For any addresses $a \in A$ and $a' \in A'$ with $f_{\mathrm{nav}}(a) =
f'_{\mathrm{nav}}(a') = b \in \Bstar$ (\ie{}, the two navigation paths
return identical bytes), $p(f_{\mathrm{nav}}(a)) = p(f'_{\mathrm{nav}}(a'))$
as values in $R$. The equality holds bit-for-bit under any canonical
serialization of $R$.
\end{theorem}

\begin{proof}
\textbf{(I).} By Definition~\ref{def:source} the parser's domain is $\Bstar$,
independent of $A$. Both $f_{\mathrm{nav}}: A \rightharpoonup \Bstar$ and
$f'_{\mathrm{nav}}: A' \rightharpoonup \Bstar$ have codomain $\Bstar$,
matching the domain of $p$. Composition in $\mathbf{PFn}$ is defined when
the codomain of the inner function matches the domain of the outer; thus
$p \circ f_{\mathrm{nav}}$ and $p \circ f'_{\mathrm{nav}}$ are both
well-typed in $\mathbf{PFn}$. The construction does not reference $\pi$ or
$\pi'$.

\smallskip\noindent
\textbf{(II).} By assumption, $p$ is deterministic and side-effect-free: $p$
is mathematically a total function on its domain modulo the partiality
inherited from $f_{\mathrm{nav}}$, with no implicit dependency on global
state. Determinism implies $p(b) = p(b)$ for any fixed $b \in \Bstar$.
Substituting $b = f_{\mathrm{nav}}(a) = f'_{\mathrm{nav}}(a')$ yields
$p(f_{\mathrm{nav}}(a)) = p(b) = p(f'_{\mathrm{nav}}(a'))$. Bit-for-bit
equality of canonical serializations follows because record canonicalization
is itself a deterministic function of the record. \qedhere
\end{proof}

Half (I) is the syntactic content -- type composition -- which has been an
implicit design principle in tools like TSK and Volatility but has not
previously been stated formally for cross-source composition. Half (II) is
the semantic content and is the load-bearing claim of the theorem: it says
that for a properly written parser, four distinct navigation paths into
the ``same'' bytes (\src{P}-live, \src{P}-image, \src{M}-carve, \src{Q}-pull
of the Chrome \texttt{History} file in Section~\ref{sec:case3}) produce
identical records, and -- because canonical serializations are
deterministic -- identical SHA-256 hashes of those records. This converts
Parser Independence from a folk principle into an empirically falsifiable
prediction.

\paragraph{Scope of the determinism assumption.} The theorem requires the
parser to be free of nondeterminism (\eg{}, hash randomization of internal
maps, system-clock-dependent fields) and free of externalized state (\eg{},
no calls to schema-resolution services). These are standard
``pure-function'' constraints; we discuss empirical compliance as a
threat-to-validity (Section~\ref{sec:discussion}).

\subsection{Layer Dependency Rule}

The dependency rule is best understood not as an independent theorem but as
the contrapositive of Theorem~\ref{thm:parser-indep}. We state it
correspondingly. The structure is: \emph{if Parser Independence is required
of a parser, then the parser must not depend on NAVIGATE.}

\begin{theorem}[Layer Dependency Rule, Contrapositive Form]
\label{thm:layer-rule}
Let $p$ be a parser, and let $\Pi$ denote the set of navigation primitives
that some deployment of $p$ may encounter. If $|\Pi| \geq 2$ and $p$ is
required to satisfy Parser Independence (Theorem~\ref{thm:parser-indep})
across $\Pi$, then $p$'s domain must equal $\Bstar$ -- equivalently, $p$
must not import or otherwise depend on any NAVIGATE component specific to a
particular $\pi \in \Pi$.
\end{theorem}

\begin{proof}
Theorem~\ref{thm:parser-indep} states that semantic invariance holds for
parsers $p : \Bstar \to R$. The hypothesis is unconditional in $\pi$: the
domain of $p$ is $\Bstar$, not any $A_\pi$. Now suppose $p$ has domain
$A_\pi$ for some particular $\pi$ -- equivalently, $p$ depends on a
NAVIGATE component fixing the address space to $A_\pi$. Then $p$ is not a
function on $\Bstar$, and Theorem~\ref{thm:parser-indep} does not apply to
$p$. We must check that $p$ also \emph{cannot} satisfy Parser Independence
under any alternative argument. Take $\pi' \neq \pi$ in $\Pi$ and a source
$S' = (A', f'_{\mathrm{nav}}, p)$ of primitive $\pi'$. The domain of $p$
contains addresses of type $A_\pi$ but not $A_\pi$-typed values produced
through $f'_{\mathrm{nav}}$, which produces values in $\Bstar$ (or in
$A'$). The composition $p \circ f'_{\mathrm{nav}}$ is therefore not even
well-typed. Hence $p$ fails Type composition (Theorem~\ref{thm:parser-indep}.I)
and \emph{a fortiori} fails Semantic invariance. The contrapositive
direction is established without circularity: we have shown that the
\emph{absence} of an unconditional-in-$\pi$ domain forces a type failure on
some pair of primitives, not that a parser ``cannot be reused'' as a
premise. \qedhere
\end{proof}

The contrapositive is the practitioner's rule of thumb: \emph{if your
parser imports a navigation crate, it cannot satisfy Parser Independence
across primitives, and you must accept that limitation explicitly}. In the
Issen implementation (Section~\ref{sec:implementation}), this is enforced
both by crate boundaries and by build-system rejection of cyclic or upward
dependencies.

\subsection{Attacker-Durability}

\begin{definition}[Attacker-Durable Evidence]
\label{def:durable}
Let $e$ be an evidential artifact captured at time $t_0$, with content
$\mathrm{content}(e)$. Let $\mathcal{A}$ be the set of actions available
to an adversary $A$. We say $e$ is \emph{attacker-durable with respect to
$A$} iff for every action $a \in \mathcal{A}$ taken at any time $t_1 > t_0$,
$\mathrm{content}(e)$ is not altered by $a$.
\end{definition}

The definition is deliberately sharp: it asks only whether the content of
the captured artifact can be altered after capture, and ignores whether the
adversary can prevent the capture, prevent its disclosure, or fabricate
alternative artifacts. (Those questions belong to a separate analysis of
\emph{capture integrity} and are out of scope here.)

For \src{P}, \src{M}, and \src{L}, durability with respect to a host-resident
adversary is established by acquisition: once the bytes are off-host, on-host
mutation cannot reach back. The interesting cases are \src{Q} and \src{C}.

\paragraph{Threat-model preliminaries for Q.} Let $\mathcal{T}_{\mathrm{host}}$
denote the set of attacker capabilities available to a host-resident
adversary: arbitrary code execution under any user (including \texttt{root}
or \texttt{SYSTEM}), modification of any disk-resident artifact,
modification of any kernel data structure, modification of any memory
page, denial of service. Let $\mathcal{T}_{\mathrm{net}}$ denote
network-resident capabilities: passive observation, active
man-in-the-middle modification of unauthenticated traffic, replay of
captured packets. Let $\mathcal{T}_{\mathrm{inv}}$ denote attacker
capabilities at the investigator side: physical access to investigator
storage, write access to that storage, key compromise of the investigator's
authentication. The Q-Durability claim is parameterized explicitly over
which subset of $\{\mathcal{T}_{\mathrm{host}}, \mathcal{T}_{\mathrm{net}},
\mathcal{T}_{\mathrm{inv}}\}$ is granted to the adversary.

\begin{proposition}[Q-Durability under bounded adversary]
\label{prop:q-durability}
Let $r$ be a result-row sequence captured by an investigator-controlled
collection at time $t_0$ via \src{Q} navigation, transmitted over channel
$\chi$ to investigator-controlled storage $\Sigma$. Suppose the adversary
possesses $\mathcal{T}_{\mathrm{host}}$ but not $\mathcal{T}_{\mathrm{net}}$
on $\chi$ (specifically: $\chi$ provides authenticated, integrity-protected
delivery -- \eg{}, mutually-authenticated TLS with a pinned certificate)
and not $\mathcal{T}_{\mathrm{inv}}$ on $\Sigma$. Then $r$ is
attacker-durable (Definition~\ref{def:durable}) against any host action at
any $t_1 > t_0$.
\end{proposition}

\begin{proof}
After $t_0$, $r$ resides on $\Sigma$. The adversary's capability set is
$\mathcal{T}_{\mathrm{host}}$, which contains no operation whose codomain
includes $\Sigma$: no host action can write to investigator-controlled
storage that is not exposed through a host-reachable interface (and
$\Sigma$ is not so exposed by hypothesis). The integrity property of
$\chi$ (assumed by hypothesis) guarantees that what $\Sigma$ received
equals what was transmitted at $t_0$. Therefore for any $a \in
\mathcal{T}_{\mathrm{host}}$ at $t_1 > t_0$, $\mathrm{content}(r)$ on
$\Sigma$ is unaffected by $a$. \qedhere
\end{proof}

\paragraph{Acknowledged hypothesis weight.}
The hypotheses ``$\chi$ provides authenticated integrity-protected
delivery'' and ``$\Sigma$ is not exposed to $\mathcal{T}_{\mathrm{inv}}$''
are doing real work, and we acknowledge their weight explicitly. The
proposition is a \emph{conditional} durability claim, not an absolute one.
If the adversary additionally possesses $\mathcal{T}_{\mathrm{net}}$ on
$\chi$ at the moment of transmission, the rows received may differ from
the rows transmitted; the claim degrades to authenticity-of-delivery
rather than authenticity-of-source. If the adversary possesses
$\mathcal{T}_{\mathrm{inv}}$ on $\Sigma$, no source type achieves
attacker-durability -- not \src{P}, \src{M}, \src{L}, \src{Q}, or even
\src{C} (the adversary can simply destroy or substitute the artifact at
the investigator's end). The boundary of the claim is therefore: \src{Q}
durability is meaningful precisely when the adversary's foothold is
host-local and the investigator's collection-and-storage infrastructure
is not part of the adversary's attack surface. This is the operational
regime in which live response is deployed; the proposition formalizes its
soundness.

The proposition formalizes a property that practitioners have understood
informally for years~\cite{cohen2011distributed}: it is precisely why a
Velociraptor sweep at the moment of incident detection is preferred over a
post-hoc disk image when the adversary still has a foothold. The disk image
captures whatever the adversary leaves behind; the live query captures the
state \emph{at} the moment of capture, which the adversary can no longer
alter.

\begin{proposition}[C-Intrinsic-Durability]
\label{prop:c-durability}
Let $H : \Bstar \to \{0,1\}^n$ be a hash function and let $\epsilon_H$
denote its collision-resistance bound, \ie{}, the maximum probability that
a polynomial-time adversary outputs $b \neq b'$ with $H(b) = H(b')$
(see~\cite{rogaway2004hash}). Let $b$ be a blob with $h = H(b)$ stored in a
content-addressed store. Then the artifact addressed by $h$ is
attacker-durable (Definition~\ref{def:durable}) with probability at least
$1 - \epsilon_H$: any adversarial modification yielding $b'$ with $h =
H(b')$ has probability at most $\epsilon_H$, and otherwise $H(b') \neq h$
and the address $h$ continues to resolve to $b$.
\end{proposition}

\begin{proof}
By the second-preimage-resistance reduction in~\cite{rogaway2004hash},
producing $b' \neq b$ with $H(b') = H(b) = h$ has probability at most
$\epsilon_H$ for any polynomial-time adversary. Conditioning on the
non-collision event (probability $\geq 1 - \epsilon_H$), the address $h$
of any correctly-implemented CAS lookup resolves only to a blob whose hash
equals $h$, which under non-collision is $b$ alone. Therefore the artifact
at $h$ remains $b$ with probability $\geq 1 - \epsilon_H$. \qedhere
\end{proof}

\paragraph{Hash-agnosticism of \src{C}.}
The proposition is parameterized in $\epsilon_H$ alone; it is not
specialized to any particular hash function. SHA-1
($\epsilon_H$ no longer negligible against chosen-prefix
attacks~\cite{stevens2017sha1,leurent2020sha1chosen}), SHA-256 (current
default for OCI~\cite{ociimagespec_digest} and Sigstore), and the BLAKE3
or post-quantum candidates likely to follow are all admitted. Git's
ongoing migration to SHA-256~\cite{git2020sha256} is therefore
accommodated within the same proposition by re-instantiating $\epsilon_H$.
A practical implementation must reject or downgrade trust for stores using
collision-broken hashes; the formal statement is robust because
$\epsilon_H$ is exposed as a parameter rather than baked in.

\begin{corollary}[\src{C} as Trust-Anchor-Free]
\label{cor:c-anchor-free}
\src{C} is the only source type whose attacker-durability requires no
external trust anchor. \src{P}, \src{M}, and \src{L} require a trusted
acquisition act and chain of custody. \src{Q} requires a trusted transmission
channel and trusted investigator storage. \src{C} requires only the
collision-resistance of the hash function, which is internal to the
addressing scheme itself.
\end{corollary}

\begin{proof}[Proof sketch]
The durability claims for \src{P}, \src{M}, \src{L} reference an
acquisition-and-custody process that is external to the bytes themselves:
without that process, the on-host adversary can mutate the source. The
durability claim of \src{Q} (Proposition~\ref{prop:q-durability}) explicitly
references the channel $\chi$ and investigator storage as trust anchors.
The durability claim of \src{C} (Proposition~\ref{prop:c-durability})
references only the hash function's collision-resistance, which is a
property of the addressing scheme. \qedhere
\end{proof}

\subsection{Hash-Pivot}

The intrinsic durability of \src{C} enables a forensic operation that has
no analog in the other primitives. The naive definition would take a union
of query results from heterogeneous stores; this is type-unsound, because
each store returns metadata under its own schema (a Sigstore Rekor entry
is not type-equal to an OCI manifest reference, which is not type-equal to
a git commit reference). We give a schema-conscious definition.

\begin{definition}[Hash-Pivot, schema-conscious]
\label{def:hash-pivot}
Let $\mathcal{C} = \{(C_1, \tau_1), (C_2, \tau_2), \ldots, (C_k, \tau_k)\}$
be a finite collection of content-addressed stores together with their
result-record \emph{schema tags}: each $\tau_i$ is a type label naming the
metadata schema that $C_i$ returns (\eg{},
\texttt{ociManifestRef}, \texttt{rekorEntry}, \texttt{gitCommitRef},
\texttt{intotoLink}, \texttt{ipfsLink}). Let $M_{\tau_i}$ denote the typed
record space associated with schema $\tau_i$. The
\emph{schema-tagged query} is
\[
\mathrm{query}(C_i,\,h) \;:\; \{0,1\}^{n} \rightharpoonup \{(b, m, \tau_i)
\mid b \in \Bstar,\, H(b) = h,\, m \in M_{\tau_i}\}.
\]
The \emph{hash-pivot} operation is the typed multiset (disjoint union)
\[
\mathrm{H\text{-}Pivot}(h)\;=\;\biguplus_{(C_i,\tau_i) \in \mathcal{C}}
\mathrm{query}(C_i,\,h),
\]
where $\biguplus$ is the multiset disjoint union over the disjoint-tagged
type $\biguplus_i \big(\Bstar \times M_{\tau_i} \times \{\tau_i\}\big)$.
\end{definition}

The disjoint union is mandatory: results from different store types are
\emph{kept distinct} by their schema tag, allowing downstream consumers to
dispatch on $\tau_i$ and reason about $m$ in its native schema. A common
implementation pattern projects each tag-segregated component into a
unifying ``provenance fragment'' record at orchestration time
(Section~\ref{sec:implementation}), but this projection is application
choice, not part of the definition.

\smallskip
The operation is forensically powerful for two reasons. First, the address
$h$ commutes across stores: the same hash refers to the same content in
every CAS by C-Intrinsic-Durability. Second, each $C_i$ that returns a
result contributes a typed provenance fragment (build identity, signing
identity, commit author, in-toto link), allowing the investigator to
assemble a chain of provenance from a single starting hash without losing
schema fidelity. Section~\ref{sec:case2} demonstrates a hash-pivot
resolving a supply-chain compromise from a suspicious-binary discovery on
disk.

\subsection{Toward a Completeness Theorem}
\label{sec:completeness-formal}

A natural question is whether the five primitives are exhaustive. We frame
the question in two parts: a partial-completeness theorem that we can prove
under structural assumptions, and an open conjecture that fully would close
the question.

\begin{theorem}[Partial Completeness over Locally-Computable Oracles]
\label{thm:partial-complete}
Let $S = (A, f_{\mathrm{nav}}, f_{\mathrm{parse}})$ be an evidence source
whose navigation function $f_{\mathrm{nav}}$ is computable by a
deterministic Turing machine~\cite{sipser2012computation} accessing a
finite set of oracles, each of which exposes one of the following access
disciplines: (i) random-access reads against a finite linearly-addressed
byte array, (ii) lookup against a virtual-to-physical page-table whose
entries are encoded in a byte array of type (i), (iii) sequential-access
reads against a linearly ordered cursor stream, (iv) request-response
dispatch against an external endpoint with input \emph{query expression}
and output \emph{result-row stream}, or (v) lookup by hash against a
finite map from $\{0,1\}^n$ to byte arrays. Then $S$'s navigation primitive
$\pi$ (Definition~\ref{def:nav-primitive}) is equivalent under $\sim$ to
one of $\pi_{\mathrm{P}}, \pi_{\mathrm{M}}, \pi_{\mathrm{L}}, \pi_{\mathrm{Q}},
\pi_{\mathrm{C}}$, depending on which oracle dominates the navigation.
\end{theorem}

\begin{proof}[Proof sketch]
The five oracle disciplines are constructed precisely to encode the address
algebra (Definition~\ref{def:addr-algebra}) of, respectively,
\src{P}, \src{M}, \src{L}, \src{Q}, and \src{C}. A navigation function
restricted to oracle (i) has address-space algebra equivalent to $A_P$
under the bijection that identifies linearly-addressable bytes with block
sequences. Likewise for the other four. A navigation function using
multiple oracles can be decomposed into a primary oracle (the one whose
address space dominates the partial-function-composition lattice in
$\mathbf{PFn}$) and auxiliary oracles; the equivalence class of the
composition is determined by the primary. Cases (i)-(v) are exhaustive
under the hypothesis. \qedhere
\end{proof}

\begin{conjecture}[Full Completeness]
\label{conj:full-complete}
Every evidence source admissible to digital forensic practice -- where
``admissible'' is defined as ``capable of being captured, transmitted, and
re-examined under chain-of-custody discipline'' -- has navigation primitive
equivalent under $\sim$ to one of $\pi_{\mathrm{P}}, \pi_{\mathrm{M}},
\pi_{\mathrm{L}}, \pi_{\mathrm{Q}}, \pi_{\mathrm{C}}$.
\end{conjecture}

The partial-completeness theorem rules out evidence sources whose
navigation functions are confined to the five oracle disciplines. The
remaining gap to full completeness is the question whether forensic
admissibility forces the oracle restriction; we conjecture that it does,
but a proof would require a formalization of forensic admissibility that
we leave to future work. We list illustrative cases and their reductions
in support of the conjecture.

\paragraph{Network packet captures.} PCAP and PCAP-NG~\cite{pcapng2024}
present a stream of records keyed by capture timestamp; this reduces to
\src{L} (oracle (iii)).

\paragraph{Blockchain ledgers.} Each block's address is its hash; references
between blocks and transactions are hashes; this is \src{C} (oracle (v)),
with the genesis block as Merkle DAG root.

\paragraph{Database tables.} Static row data on a disk is reached via
\src{P}; live queries against a running database engine are \src{Q};
the underlying B-tree pages on storage are \src{P}.

\paragraph{Cloud audit logs.} CloudTrail and equivalent services expose a
stream of records keyed by event timestamp and pagination cursor; this is
\src{L}, with the optional \src{Q} layering when the investigator pulls via
the audit-API rather than from an exported store.

\paragraph{Mobile device extractions.} Decoded backups land on a host
filesystem (\src{P}) once extracted; live agent queries against an unlocked
device are \src{Q}.

\paragraph{Stream-processing materialized state.} Continuous-query systems
(Kafka Streams, Flink) maintain a running state that is neither persisted
nor produced ad-hoc by an investigator query. We classify these as a
\src{L}-input, \src{Q}-output composition: the underlying stream is \src{L}
(oracle (iii)) and the investigator's snapshot of materialized state is
\src{Q} (oracle (iv)).

\paragraph{Quantum-entangled storage.} Theoretically interesting; not
currently a practical evidence substrate; the oracle restriction in
Theorem~\ref{thm:partial-complete} is over deterministic Turing machines
and would need extension for a formal account.

The remaining open question is whether Conjecture~\ref{conj:full-complete}
holds against \emph{every} new computing paradigm: a sixth primitive, if
one exists, would be one whose address-space algebra cannot be simulated
under any of the five oracle disciplines. We discuss this further in
Section~\ref{sec:discussion}.
